% answers_short.tex —— 2012 年·辽宁卷·理科 · 简短答案速查
% 由 tools/make_short_answers.py 自 answers.tex 的 \gksolhead/\gkanswer 提取，
% 并与 2012/2012_数学真题_辽宁理_答案要点.md 逐题比对后生成，勿手改（改源 MD 后重跑）。

\begin{gknotice}
\textbf{使用说明}：本节只给最终答案便于快速核对，完整推导见解析版。答案为编写者独立推导，\textbf{非官方标准答案}。
\end{gknotice}

\gkpart{选择题}
\begin{center}\begin{tabular}{ccccccccccccc}
\toprule
题号 & 1 & 2 & 3 & 4 & 5 & 6 & 7 & 8 & 9 & 10 & 11 & 12 \\
\midrule
答案 & B & A & B & C & C & B & A & D & D & C & B & C \\
\bottomrule
\end{tabular}\end{center}

\gkpart{填空题}
\begin{itemize}[leftmargin=2.6em]
  \item[\textbf{13.}] $38$。几何体是 $4\times3\times1$ 长方体挖去一个直径 $2$ 的圆柱形通孔，外表面积 $2(12+4+3)=38$，孔壁面积与上下两个圆孔面积相抵。
  \item[\textbf{14.}] $2^n$。由 $2(1+q^2)=5q$ 得 $q=2$ 或 $\dfrac12$，递增取 $q=2$；再由 $a_5^2=a_{10}$ 得 $a_1=2$。
  \item[\textbf{15.}] $-4$。两切线为 $y=4x-8$ 与 $y=-2x-2$，交点 $A(1,-4)$。
  \item[\textbf{16.}] $\dfrac{\sqrt3}{3}$。设 $PA=PB=PC=l$，外接球半径 $\dfrac{\sqrt3}{2}l=\sqrt3$ 得 $l=2$；球心到平面 $x+y+z=2$ 的距离为 $\dfrac{|3-2|}{\sqrt3}=\dfrac{\sqrt3}{3}$。
\end{itemize}

\gkpart{解答题}
\begin{itemize}[leftmargin=2.6em]
  \item[\textbf{17.}] (1) $\dfrac12$。$2B=A+C$ 且 $A+B+C=\pi$，故 $B=\dfrac\pi3$，$\cos B=\dfrac12$；(2) $\dfrac34$。由 $b^2=ac$ 及正弦定理得 $\sin^2B=\sin A\sin C$，故 $\sin A\sin C=\dfrac34$
  \item[\textbf{18.}] (1) 见解析版；(2) $\sqrt2$。取 $AA'=1$，两半平面的法向量分别为 $(1,-1,\lambda)$、$(-3,-1,\lambda)$，由法向量垂直得 $\lambda^2=2$
  \item[\textbf{19.}] (1) 列联表为男 $28,17,45$，女 $45,10,55$，合计 $73,27,100$；$\chi^2=\dfrac{940900}{195129}\approx4.822>3.841$，有 $95\%$ 的把握认为“体育迷”与性别有关；(2) $X\sim B\left(3,\dfrac{27}{100}\right)$，分布列为 $P(X=0)=\dfrac{389017}{1000000}$，$P(X=1)=\dfrac{431649}{1000000}$，$P(X=2)=\dfrac{159651}{1000000}$，$P(X=3)=\dfrac{19683}{1000000}$；$E(X)=\dfrac{81}{100}$，$D(X)=\dfrac{5913}{10000}$
  \item[\textbf{20.}] (1) $\dfrac{x^2}{a^2}-\dfrac{y^2}{b^2}=1$（$x<-a$）。设 $A(-x_1,y_1)$，$B(-x_1,-y_1)$，联立两直线得 $M\left(-\dfrac{a^2}{x_1},-\dfrac{ay_1}{x_1}\right)$，消去 $x_1,y_1$ 即得；(2) 见解析版
  \item[\textbf{21.}] (1) $a=0$，$b=-1$。由 $f(0)=0$ 得 $b=-1$，由 $f'(0)=\dfrac32$ 得 $1+\dfrac12+a=\dfrac32$，故 $a=0$；(2) 见解析版
\end{itemize}

\gkpart{选做题}
\begin{itemize}[leftmargin=2.6em]
  \item[\textbf{22.}] (1) 见解析版；(2) 见解析版
  \item[\textbf{23.}] (1) $C_1:\rho=2$，$C_2:\rho=4\cos\theta$；交点极坐标为 $\left(2,\dfrac\pi3\right)$，$\left(2,-\dfrac\pi3\right)$；(2) $\begin{cases}x=1,\\y=t\end{cases}$（$-\sqrt3\leqslant t\leqslant\sqrt3$）。两圆方程相减得公共弦所在直线 $x=1$，再与 $x^2+y^2=4$ 联立
  \item[\textbf{24.}] (1) $a=2$。$|ax+1|\leqslant3$ 即 $-4\leqslant ax\leqslant2$，与 $-2\leqslant x\leqslant1$ 比较得 $a=2$；(2) $k\geqslant1$。令 $h(x)=|2x+1|-2|x+1|$，分段求得 $h(x)\in[-1,1]$，故 $k$ 的取值范围是 $[1,+\infty)$
\end{itemize}
